\documentclass[10pt,a4paper]{amsart}
\usepackage[margin=19mm]{geometry}
\usepackage{amsmath,amssymb,mathtools}
\usepackage[scr=rsfs,cal=euler]{mathalfa}
\usepackage{xfrac}
\usepackage[unicode,hidelinks,bookmarks=false]{hyperref}
\newcommand{\ie}{i.e.\ }
\newcommand{\cf}{cf.\ }
\newcommand{\resp}{respectively\ }
\newcommand{\Subsection}{\subsection}
\providecommand{\nobreakdash}{\nobreak\hbox{-}}
\setlength{\emergencystretch}{2em}
\multlinegap0pt
\usepackage{bm}
\usepackage{bbm}
\usepackage{dsfont}
\usepackage{xspace}
\usepackage{cancel}
\usepackage{mathtools}
\usepackage{amsthm}
\makeatletter
\@ifundefined{proposition}{\newtheorem{proposition}{Proposition}}{}
\makeatother
\newcommand{\bI}{\mathbf{I}}
\newcommand{\bM}{\mathbf{M}}
\title[Parallel Stochastic Newton Method]{Parallel Stochastic Newton Method}
\author{Mojm\'ir Mutn\'y, Peter Richt\'arik}
\date{Source: arXiv:1705.02005v3 (CC BY 4.0)}
\begin{document}
\maketitle

\noindent\emph{Source and licence.} Parallel Stochastic Newton Method. Version arXiv:1705.02005v3 (CC BY 4.0), distributed under the Creative Commons Attribution 4.0 International licence, as recorded in the arXiv licence metadata for this exact version and shown on the abstract page https://arxiv.org/abs/1705.02005v3. Full rights notice: licenses/arxiv-1705.02005-CC-BY-4.0.txt.\par\smallskip

\noindent\textit{Editorial scope.} Counted unit: the Proposition whose source label is \texttt{None} in the article (source0.tex lines 604--621, source label \texttt{None}), delivered together with the article's own complete original proof of it (source0.tex lines 622--661, one \texttt{proof} environment of 40 lines). No mathematics is written, repaired or re-derived: statement and proof are the article's own text line for line. Required context retained uncounted: none. Every definition, display and label of those lines is retained unchanged. Omitted from this package and not redistributed: 1--603, 662--1025. Nothing inside a retained excerpt is omitted or altered: every retained body is delivered exactly as the article's lines are written, so no omission receipt of any kind accompanies this package. Citations: the retained excerpts carry no citation command at all, so no bibliography entry of the article is transcribed and no reference is redistributed. Reference adaptation: none was needed and none was made; every reference inside the delivered unit resolves inside it, so no unresolved reference or citation is produced. Printed slips kept literal: no printed word, symbol, hypothesis or number of the counted statement or of its proof was altered. Layout and class: the article's own class is \texttt{article} with packages including parskip, color, inputenc, fontenc, hyperref, url, booktabs, amsfonts, nicefrac, microtype, subcaption, algorithm, algpseudocode, braket; the wrapper is the lane's pinned \texttt{amsart} editorial wrapper at 10pt on A4, which restates the theorem-like environments the retained text needs and adds the article's own macro definitions that the retained text needs (\texttt{\textbackslash bI}, \texttt{\textbackslash bM}), with extra wrapper packages (bm, bbm, dsfont, xspace, cancel, mathtools). The delivered numbering is the unit's own. Assets: no retained span contains a figure environment, a graphics-inclusion command, a PDF-page inclusion, a table or a TikZ picture; no image file is copied into this package, the package declares no asset dependency and no image file is redistributed with it. Licence: version arXiv:1705.02005v3 of this article is distributed under the Creative Commons Attribution 4.0 International licence, as recorded in the arXiv licence metadata for this exact version and shown on the abstract page https://arxiv.org/abs/1705.02005v3. A full rights notice is delivered at licenses/arxiv-1705.02005-CC-BY-4.0.txt. Characters: every retained excerpt is pure ASCII, so the delivered engine, XeLaTeX with its default Latin Modern text font, reports no missing character.
	\begin{proposition}
		Suppose we have a $\rho$-matrix $\mathbf{M}$ with $0<\rho<1$ as a parameter. Then $\sigma_1$ and $ \theta$ can be expressed as functions of $n$ and $\tau$ for $\tau$-nice sampling as follows 
		\begin{equation}
			\sigma_1=(1-\rho)\left(1-\frac{(\tau-1)B(\tau)}{(n-1)A(\tau)}\right)\frac{\tau A(\tau)}{n}
		\end{equation}
		and 
		\begin{equation}
			 \theta=(n\rho-\rho+1)\left(n\frac{(\tau-1)B(\tau)}{(n-1)A(\tau)}-\frac{(\tau-1)B(\tau)}{(n-1)A(\tau)}+1\right)\frac{\tau A(\tau)}{n}
		\end{equation}
		where we use function
		\begin{equation}
			A(\tau)=\frac{((\tau-2)\rho+1)}{(1-\rho)((\tau-1)\rho+1)}
		\end{equation}
		and 
		\begin{equation}
			B(\tau)=-\frac{\rho}{(1-\rho)((\tau-1)\rho+1)}.
		\end{equation}
	\end{proposition}

	\begin{proof}
		There are $n$ eigenvalues of the matrix $\mathbf{M}$; $n-1$ of them are degenerate and equal to $1-\rho$. This can be verified by plugging the suitable eigenvector which has only two non-zero elements ${-1,1}$. The last eigenvalue can be computed using the trace of the matrix. $\lambda=n-(1-\rho)(n-1)$. After solving for this eigenvalue, we immediately see it is the biggest one, thus, $\lambda_{\max}(\mathbf{M})=n\rho-\rho+1$.
		
		Given a $\tau$-nice sampling with $\tau$ as parameter, the matrix that is sub-sampled is a $\rho$-matrix of the size $\tau \times \tau$. Due to the symmetry of the problem, any subset of $\tau$ coordinates slices out the same matrix. Additionally, the cofactor matrices for diagonal elements are all equal, and the cofactor matrices for off-diagonal elements are all equal as well.
		
		As the entries in the inverse of the matrix depend on the determinant of the cofactor matrix only, all diagonal and non-diagonal elements of the inverse are equal. In addition, taking expectation does preserve this structure as the $\tau$-nice sampling is uniform. Thus, Let us denote the value of these two entries in diagonal and off-diagonal in the sampled inverted matrix as $A(\tau)$ and $B(\tau)$ respectively.
		
		Therefore, the value of $\mathbb{E}[\mathbf{M}_{\hat{S}}^{-1}]=\frac{\tau}{n}\bI	A(\tau)+\frac{\tau(\tau-1)}{n(n-1)}(\mathbf{S}-\bI)B(\tau)$, where $\mathbf{S}$ is matrix full of ones. To determine the values of $A$ and $B$, we compute inverse for a given $\tau$ using Cramer's rule (using matrix minors)
		
		\begin{equation}\label{eq:cramer}
			\bM^{-1}=\frac{1}{\det(\bM)}\mathbf{C}^\top,
		\end{equation}
		where $\mathbf{C}$ is matrix of cofactors. 
		
		The trick here is that we need to look only at two cofactors due to the structure of the matrix. So for $A(\tau)$ on diagonal we take $\mathbf{C}_{11}$ minor. We adopt notation where we use subscript of $\mathbf{M}$ to denote the size of the sub-matrix of $\rho$-matrix $\mathbf{M}$. We calculate determinant as a product of the known eigenvalues discussed earlier. 
		\begin{equation*}
			A(\tau)=(\mathbf{M}_{\tau}^{-1})_{11}\stackrel{\eqref{eq:cramer}}=\frac{\det(\mathbf{M}_{\tau-1})}{\det(\mathbf{M}_{\tau})}=\frac{(1-\rho)^{(\tau-2)}((\tau-2)\rho+1)}{(1-\rho)^{(\tau-1)}((\tau-1)\rho+1)}=\frac{((\tau-2)\rho+1)}{(1-\rho)((\tau-1)\rho+1)} 
		\end{equation*}
		
		When we look at $B(\tau)$ the minor of the matrix that determines any off-diagonal is a matrix which contains only $\rho$'s at first row and continues with other rows as in matrix $\mathbf{M}_\tau$. Let us denote such type of matrix $\mathbf{N}_{\tau}$, where the subscript $\tau$ symbolizes the size of the matrix again. For $\det(\mathbf{N}_{\tau})$, we arrive at the following recurrence relation by using Laplace decomposition of determinants.
		\begin{eqnarray} \label{eq:reccur}
			\det(\mathbf{N}_{\tau})& =& (1-\rho)\mathbf{N}_{\tau-1} \\
			\det(\mathbf{N}_{2})& = & \rho^2-\rho
		\end{eqnarray}
		Therefore,
		\begin{eqnarray}
			B(\tau)&=&(\mathbf{M}_{\tau}^{-1})_{12}=-\frac{\det(\mathbf{N}_{\tau})}{\det(\mathbf{M}_{\tau})}\stackrel{\eqref{eq:reccur}}=-\frac{(1-\rho)^{(\tau-2)}\rho}{(1-\rho)^{(\tau-1)}((\tau-1)\rho+1)} \notag \\&= &-\frac{\rho}{(1-\rho)((\tau-1)\rho+1)} 
		\end{eqnarray}
		
		The matrix $\mathbb{E}[\mathbf{M}_{\hat{S}}^{-1}]$ can be manipulated to the from of $\rho$-matrix by dividing by the factor $\frac{\tau A(\tau)}{n}$. Thus, we can define a new $\rho_N$ which serves as parameter of the new $\rho$-matrix. 
		
		This parameter allows us to compute the eigenvalues of the matrix, and thus, also maximal and minimal eigenvalues. From the $\sigma_1$ definition, $\sigma_1$  is the smallest eigenvalues of a product of two $\rho$-matrices one with $\rho$ as parameter and one with $\rho_N$ as the parameter. The product of two $\rho$-matrices is a $\rho$-matrix again, and as consequence of this, the smallest eigenvalue of the product is just a product of the two smallest eigenvalues ones.
		\begin{equation}
			\sigma_1=(1-\rho)(1-\rho_N)\frac{\tau A(\tau)}{n}=(1-\rho)\left(1-\frac{(\tau-1)B(\tau)}{(n-1)A(\tau)}\right)\frac{\tau A(\tau)}{n}.
		\end{equation}
		And similarly $ \theta$, and its definition imply that $ \theta$ is product of the two biggest eigenvalues
		\begin{eqnarray}
			 \theta&=&(n\rho-\rho+1)(n\rho_N-\rho_N+1)\frac{\tau A(\tau)}{n} \notag\\&=&(n\rho-\rho+1)\left(n\frac{(\tau-1)B(\tau)}{(n-1)A(\tau)}-\frac{(\tau-1)B(\tau)}{(n-1)A(\tau)}+1\right)\frac{\tau A(\tau)}{n}	.
		\end{eqnarray}	
	\end{proof}

\end{document}
